\centerline{\bf \Huge Основы ТЧ}

\begin{flushright}
        \scriptsize{Эпизод 2. Незнакомые потолки}
\end{flushright}

\textbf{\large Теоретический блок}

\problem{1}
\textbf{Малая теорема Ферма (двумя разными способами без теоремы Эйлера).}
Если $p-$ простое число и $a-$ целое число, не делящееся на $p$, то $a^{p-1}-1$ делится на $p$.

\problem{2}
\textbf{Теорема Эйлера.}
Если $a$ и $m$ взаимно просты, то $a^{\varphi(m)} \equiv 1\pmod{m}$, где $\varphi(m)$ - функция Эйлера.

\problem{3}
\textbf{Теорема Вильсона.}
$(p-1)!\equiv-1\pmod{p}$ тогда и только тогда, когда $p$ --- простое.

\problem{4}
\textbf{Тождество Гаусса.}
$\sum_{d \mid n} \varphi(d)=n$.

\problem{5}
\textbf{Усиленная теорема Эйлера.}
Положим $L(m)=\operatorname{HOK}\left(\varphi\left(p_1^{\alpha_1}\right), \ldots, \varphi\left(p_k^{\alpha_k}\right)\right)$. Тогда, если $(a, m)=1$, то $a^{L(m)}-1$ делится на $m$.

\problem{6}
\textbf{Усиленная теорема Вильсона.}
$(p-m)!\cdot(m-1)!\equiv(-1)^m \pmod{p}$

\textbf{\large Практический блок}

\problem{1}
Найдите все пары чисел $(p, n)$, где $p$ - простое, а $n$ - целое и при этом $p^4+211=n^2+5 p n$. В ответе укажите значения $n$. Если их несколько, перечислите их в любом порядке через запятую.

\problem{2}
Натуральные числа вида $\underbrace{11 \ldots 1}_n$ (десятичная запись состоит из $n$ единиц) будем обозначать $R_n$. Докажите, что существует такое натуральное число $k$, что $R_n$ делится на 41 тогда и только тогда, когда $n$ делится на $k$.

\problem{3}
На экране суперкомпьютера напечатано число $11 \ldots 1$ (900 единиц). Каждую секунду суперкомпьютер заменяет его по следующему правилу. Число записывается в виде $\overline{A B}$, где $B$ состоит из двух его последних цифр, и заменяется на $2 \cdot A+8 \cdot B$ (если $B$ начинается на нуль, то он при вычислении опускается). Например, 305 заменяется на $2 \cdot 3+8 \cdot 5=46$. Если на экране остаётся число, меньшее 100 , то процесс останавливается. Правда ли, что он остановится?

\newpage

\hspace{3mm}

\problem{4}
Про простые числа $p$ и $q$ известно, что

$$
p+p^2+\ldots p^q=q+q^2+\ldots q^p
$$


Докажите, что $p=q$.

\problem{5}
Найдите все тройки взаимно простых в совокупности натуральных чисел $(a, b, c)$ таких, что для любого натурального $n$ число $\left(a^n+b^n+c^n\right)^2$ делится на $a b+b c+c a$.

\problem{6}
(a)
\textit{Лемма.}
Пусть $\varphi(n)$ --- функция Эйлера числа $n$. (Количество чисел от 1 до $n$ взаимно простых с $n$ ). Тогда для любого натурального числа $n>1$ справедливо неравенство

$$
\sum_{d \mid n} d<\frac{n^2}{\varphi(n)} .
$$

\noindent
(b)
Сумма всех натуральных делителей числа $n$ более чем в 100 раз превосходит само число $n$. Докажите, что есть сто идущих подряд чисел, каждое из которых имеет общий делитель с $n$ больший 1.

\problem{7}
В строку выписаны 200 натуральных чисел. Среди любых двух соседних чисел строки правое либо в 9 раз больше левого, либо в 2 раза меньше левого. Может ли сумма всех этих 200 чисел равняться $24^{2022}$?